\chapter{Introduction}
\section{The reporting problem}

An investor holding stakes in a dozen young companies has to say, in writing and on a date, how each
one is doing. The awkward part is not writing the summary but standing behind it: a figure that came
from a slide deck, a funding round that was announced but never completed, or an accounts filing that
had not yet been made can each turn a confident sentence into one that cannot be defended when someone
asks where it came from. The difficulty is structural rather than careless, because early-stage
records are themselves incomplete and delayed \citep{hardman2023small}, and a statement can carry a
citation and still not be supported by it \citep{gao2023alce}. That question -- where did this claim
come from, and does the source actually say it -- is the problem this dissertation addresses.

The verified project charter describes three reporting activities: someone gathers company
information, someone turns it into an executive summary, and someone decides whether that summary
may be released. Table~\ref{tab:intro-business-process} sets out who does what. These are jobs to be
done rather than assumed job titles. The legal company and the reporting date must both be clear,
because being incorporated does not prove that a business is trading, and recent records may be
missing simply because a filing is not yet due \citep{galanakis2026chrt,hardman2023small}.

\begin{table}[htbp]
  \centering
  \small
  \setlength{\tabcolsep}{4pt}
  \renewcommand{\arraystretch}{1.12}
  \begin{tabularx}{\textwidth}{@{}P{0.15\textwidth}P{0.17\textwidth}P{0.19\textwidth}P{0.16\textwidth}Y@{}}
    \toprule
    \textbf{Actor} & \textbf{Input} & \textbf{Activity and hand-off} & \textbf{Output} & \textbf{Pain point addressed} \\
    \midrule
    Data contributor / reporting operator & Company submission and eligible public records & Collect information and hand it to aggregation & Source records and typed observations & Mixed formats, uncertain identity, blanks and source gaps \\
    Reporting operator & Typed observations and evidence & Aggregate and analyse, then hand a draft to review & Executive-level draft with claim links & Period mismatch, unsupported claims and hidden conflicts \\
    Reviewer / approver & Draft, evidence, conflicts and gaps & Approve, reject or return the current version & Approved report for the downstream user & Unclear release authority and unrecorded corrections \\
    \bottomrule
  \end{tabularx}

  \caption{Verified portfolio-reporting process and main control point}
  \label{tab:intro-business-process}
  \vspace{3pt}

  \parbox{\textwidth}{\scriptsize
    \textit{Source note:} Actor and process summary from \texttt{docs/PROJECT\_CHARTER.md}. The
    repository evidence does not establish reporting frequency, labour cost, error rate or job title.}
\end{table}

Portfolio reports should record funding and public support as dated events with clear sources. In the
Crowdcube study, an earlier successful round may act as a signal, but its meaning still depends on the
campaign and platform \citep{wasti2024successive}. UK research on public support found later
differences between participating firms, but it also noted selection effects and did not claim that
the support caused those differences \citep{mondal2025state,thorne2026funding}. A round, grant or
award can support a dated fact. It cannot prove future success or provide a complete funding history.

Evidence about a company's development comes from dated events rather than one complete profile.
Crowdfunding results vary by pitch, platform, country, sector, region and time. Published samples may
also include only successful pitches \citep{estrin2024digital}. Research on repeat funding uses the
campaign round as its unit \citep{wasti2024successive}. UK funding data can connect opportunities,
applications, decisions, organisations, projects and outcomes, but its coverage and links remain
incomplete \citep{thorne2026funding}. A report should keep each event tied to its date and source.

These limits show that reliable reporting starts with the data. A company register may be suitable
for one task but not another \citep{krasikov2020ready}. Even required Companies House fields can be
empty, so quality rules must be linked to the intended use and checked in code
\citep{nikiforova2020quality}. The system must record the source, date, location and check result for
each important statement, and show when evidence is missing or unclear.

The main data problem is to match different records to the correct legal company and to use only
information that was available at the relevant date. Research that joined several registers used
normalised names, approximate matching and manual review, but some records still could not be matched
\citep{surak2026gateways}. Newer companies may also have no accounts because the filing date has not
arrived \citep{hardman2023small}. The prototype waits until the company number and legal identity are
clear. It keeps different types of missing data instead of changing blanks to zeros.

Text can sound convincing and still be wrong, even when it includes citations. A citation marker
alone does not guarantee factual correctness, citation precision or passage support
\citep{gao2023alce}. A statement may also match a source that is itself wrong \citep{gao2023rarr}.
Public web pages can contain hidden instructions that try to change a model's summary, source choice
or tool actions \citep{greshake2023indirect}. Web search is therefore used only to find possible
sources. Each report claim must point to a saved source and an exact passage, and a separate check
must confirm the support.

The same three properties -- factual correctness, citation precision and passage support -- are why
a report cannot treat a fluent summary as evidence. Retrieval-augmented generation can extend search
and still return irrelevant or hostile passages, and a citation marker on the resulting sentence
does not repair that \citep{gao2023ragsurvey,gao2023alce}. The prototype therefore keeps finding,
checking, extracting and verifying as separate steps.

Human approval creates a release boundary, but it does not guarantee a better result. General
guidance recommends explaining what the system can do, showing uncertainty and allowing people to
correct or dismiss its output \citep{amershi2019guidelines}. In one controlled task, some participants
followed incorrect AI advice \citep{bucinca2021forcing}. The prototype shows conflicts, missing
evidence and exact source locations, and it requires a named reviewer before export. Whether review
improves outcomes remains a future pilot question
\citep{amershi2019guidelines,bucinca2021forcing}.

Those three problems -- identity, unsupported text, and over-trusted review -- are why the aim is an
evidence-first workflow rather than a better writer. A system that cannot refuse claim admission can
only produce more sentences. A system that can refuse them can still be wrong, but the wrongness is
then a recorded decision rather than a fluent gap \citep{gao2023alce,nist2023airmf}. Chapter~3 turns
that refusal into a measurable C1/C2 difference under matched inputs, and Chapter~5 reports it on
designed cases rather than on a sample of companies \citep{ribeiro2020checklist,fu2026agents}. That
equal-input comparison is the warrant taken from the 2026 preprint, not a claim that extra roles are
generally better.

\section{Why it matters}

The reporting problem matters because an untraced claim consumes monitoring attention without leaving
a defensible record. Section~\ref{sec:reporting-problem} locates that cost in the UK equity market
and in weekly portfolio work; the figures themselves are stated only there
\citep{britishbusinessbank2025equity,gompers2016vcdecisions}.

\section{Aim, research questions and objectives}
\label{sec:research-questions}

The aim of this dissertation is to design, implement and evaluate an evidence-first, role-separated
workflow for early-stage portfolio reporting. The study examines how the prototype preserves evidence
from intake to a reviewable report and how a separate verification stage changes claim admission in
controlled cases. This follows design-science research, which links a practical problem to the
development and evaluation of an artefact \citep{hevner2004design,peffers2007dsrm}.

\subsection*{Primary research question}

\textit{How can an evidence-first, role-separated workflow produce traceable early-stage portfolio
reporting claims, and what effect does a separate verification stage have on claim admission in
controlled cases?} The question concerns the implemented workflow rather than one model, agent label
or interface. It is answered by combining RQ1 and RQ2 with the executed engineering and D0 evidence,
as required in design-science evaluation \citep{hevner2004design,peffers2007dsrm}.

\subsection*{RQ1: Input-to-claim accuracy and reliability}

\textit{How does the implemented workflow preserve company identity, reporting time, missing states,
source provenance and approval boundaries from input to report?} RQ1 is answered with implementation
records, engineering tests and labelled D0 cases. The assessment covers exact identity holds, value
and period semantics, source-record completeness, claim-to-evidence links and repeated outcomes
\citep{nikiforova2020quality,gao2023alce}.

\subsection*{RQ2: Effect of separate verification}

\textit{What effect does a separately recorded verification stage have on supported-claim precision,
unsupported claims, conflict handling, source-record completeness and repeatability in D0?} RQ2 uses
a paired comparison between C1, which has no separate verifier, and C2, which adds verification.
Both conditions receive the same fictional cases and proposed values
\citep{gao2023alce,gao2023rarr}.

\subsection*{Scope narrowing}

The original contract also proposed manual C0, named-review C3, live-web and managed-platform
comparisons. They require authorised participants, frozen references, live access and recorded time.
Those observations were not available, so these comparisons are not active research questions
\citep{amershi2019guidelines,bucinca2021forcing}. Section~\ref{sec:limitations} is the only place
those absences are discussed at length.

\subsection*{Research objectives and evaluation plan}

Four objectives connect the questions to the study. O1 defines shared rules for legal identity,
dates, missing data and source records. O2 implements those rules in a workflow that stops when
required evidence is missing or invalid. O3 evaluates the implemented controls and the C1/C2 verifier
difference using the fixed D0 fixture. O4 distinguishes supported, unsupported, conflicting,
unavailable and missing evidence in every result
\citep{hevner2004design,peffers2007dsrm}. Why the workflow separates production from judgement is
argued in Section~\ref{sec:design-alternatives}; that argument is not repeated here.

Table~\ref{tab:research-contract} summarises this evaluation plan and the evidence currently available.

\begin{table}[htbp]
  \centering
  \caption{Research contract and current evidence boundary}
  \label{tab:research-contract}
  \footnotesize
  \setlength{\tabcolsep}{2.5pt}
  \renewcommand{\arraystretch}{1.22}
  \hyphenpenalty=4000
  \exhyphenpenalty=4000
  \begin{tabularx}{\textwidth}{@{}P{0.095\textwidth}P{0.19\textwidth}Y P{0.16\textwidth}P{0.19\textwidth}@{}}
    \toprule
    \rowcolor{WMGBlue!10}
    \textbf{Question} &
    \textbf{Evaluand / comparison} &
    \textbf{Observable measures} &
    \textbf{Current evidence status} &
    \textbf{Exclusion / interpretive boundary} \\
    \midrule
    Primary question &
    C0 manual; C1 deterministic, no independent verifier; C2 bounded independent verification; C3 is C2 plus named review. &
    Automation extent; quality; provenance; reliability; failures; active and elapsed time. &
    D0: C1/C2 mechanism evidence. Human C0/C3 and D1/D2: \textbf{EVIDENCE HELD / NO-GO}. &
    Manual comparison requires authorised C0/C3. Live public-web company-research evaluation is outside RQ1--RQ3; no investment, production or exhaustive-web inference. \\
    \addlinespace
    \rowcolor{DraftPale}
    RQ1 &
    Accurate and reliable bounded transformation of heterogeneous inputs against a frozen reference. &
    Type, extraction, normalisation, missing-state, identity and time accuracy; provenance and locator completeness; outcome-vector repeat consistency. &
    D0 labelled mechanism evidence. D1 pilot and sealed D2: \textbf{HELD}. &
    D0 supports mechanism claims only; no real-portfolio, live-source or production inference. \\
    \addlinespace
    RQ2 &
    C1 without a verifier versus C2 bounded independent verification, under the same information boundary. &
    Supported precision and recall; unsupported rate; conflicts; abstention; provenance completeness; repeat consistency. &
    C1/C2 executable on labelled D0 cases. D1/D2: \textbf{HELD}. &
    Estimate verifier effect; do not assume benefit. Report trade-offs, failures and null outcomes. \\
    \addlinespace
    \rowcolor{DraftPale}
    RQ3 &
    Compare C1/C2 automation and C3 named review with manual C0. Sub-comparison: C2 before versus C3 after review. &
    Correctness and completeness; active/elapsed time; edits; approval time; reviewer utility. &
    D0 C0/C3 measures are null. Human C0/C3 and D1/D2: \textbf{EVIDENCE HELD / NO-GO}. &
    Requires ethics/data authority, frozen reference and authorised reviewers; no retrospective timing or invented participant evidence. \\
    \bottomrule
  \end{tabularx}
  \vspace{3pt}

  \parbox{\textwidth}{\scriptsize
    \textit{Interpretive boundary:} Live public-web company-research evaluation is outside RQ1--RQ3. Investment outcomes, production readiness and exhaustive-web coverage are not evaluands.\\
    \textit{Source note:} Author synthesis of frozen project contracts, not results. Evidence-status labels distinguish executable D0 mechanism evidence from held D1/D2, manual and human evidence.
  }
\end{table}
\clearpage


\section{Scope and boundary}

The system is a local, single-user research prototype rather than a deployed service. Restricted or
internal material stays local and is never sent to public generative-AI tools
\citep{cddo2023genai,nist2023airmf}. The project excludes buy or sell recommendations, speculative
valuations, profiles of people, automatic publication and claims of production readiness
\citep{fca2024promotions,greshake2023indirect}. Those boundaries are stated once;
Section~\ref{sec:limitations} is the only later limitations section.

\section{Contribution and roadmap}

The contribution has three parts: an evidence-first prototype, an executed controlled evaluation, and
a reporting method that traces claims to sources \citep{hevner2004design,peffers2007dsrm}. The
dissertation title names the implemented role structure. It is not a claim that extra agents
outperform a single program. Chapter~2 reviews the literature and argues the design alternatives,
including matched-input evidence for role separation. Chapter~3 defines the executed evaluation, the
evidence-scope table and the ethics boundary. Chapter~4 retells the artefact as the life of a claim.
Chapter~5 reports engineering and D0 results. Chapter~6 interprets RQ1 and RQ2, specifies a
prospective business pilot, and states the study's limits. Chapter~7 closes the argument. This order
follows design science by linking the problem, artefact, evaluation and communication
\citep{hevner2004design,peffers2007dsrm}.
